% Chapter 5 exercises. Source PDF pages 163--166 / printed pages 151--154.
\MathBookSourcePage{163}
\subsection*{习题}
\label{sec:ch05-exercises}
\setcounter{exerciseitem}{0}

\begin{MBExercise}
\label{exr:ch05-01}
证明 $\widetilde L^2(\Omega):=\{\phi\in L^2(\Omega):\int_\Omega\phi\,dx=0\}$ 中任意函数都可写成 $\widetilde{\mathcal D}(\Omega):=\{\phi\in\mathcal D(\Omega):\int_\Omega\phi\,dx=0\}$ 中函数的极限. 提示: 取 $\mathcal D(\Omega)$ 中一个在 $L^2$ 中收敛的序列, 再修改它以得到平均值为零的序列.
\end{MBExercise}

\begin{MBExercise}
\label{exr:ch05-02}
设 $f,g\in H^1(\Omega)$. 证明对所有 $|\alpha|=1$, 弱导数 $D_w^\alpha(fg)$ 存在, 且
\[
\MathBookFormulaID{formula-ch05-exercise-02-product-rule}
D_w^\alpha(fg)=D_w^\alpha f\,g+fD_w^\alpha g.
\]
提示: 先对 $g\in C^\infty(\overline\Omega)\cap H^1(\Omega)$ 证明, 再使用稠密性结果~\ref{thm:ch01-meyers-serrin-density}.
\end{MBExercise}

\begin{MBExercise}
\label{exr:ch05-03}
陈述并证明命题~\ref{prop:ch05-sobolev-integration-by-parts} 和~\ref{prop:ch05-green-identity} 的适当 $L^p$ 版本, 其中 $p\ne2$.
\end{MBExercise}

\begin{MBExercise}
\label{exr:ch05-04}
证明若 $u\in W_p^k(\Omega)$ 且 $|\alpha|\leq k$, 则 $D^\alpha u\in W_p^{k-|\alpha|}(\Omega)$. 再证明若 $|\alpha|+|\beta|\leq k$, 则 $D^\beta D^\alpha u=D^\alpha D^\beta u$.
\end{MBExercise}

\MathBookSourcePage{164}
\begin{MBExercise}
\label{exr:ch05-05}
为求解势流方程 $\Delta u=0$, E. Wu 在 ``A cubic triangular element with local continuity---an application in potential flow'', Int. J. Num. Meth. Eng. 17 (1981), 1147--1159 中提出如下方法. 对一个单元, 取 $K$ 为三角形, $P$ 为满足
\[
\MathBookFormulaID{formula-ch05-exercise-05-zero-net-flux-polynomial-space}
\int_{\partial K}\frac{\partial v}{\partial\nu}\,ds=0
\]
的三次多项式 $v$ 的空间, $N$ 为在每个顶点对 $v$ 及其梯度的求值. 注意对 $v\in P$, 向量场 $\nabla v$ 通过 $K$ 的净通量为零, 即流入的总量必等于流出的总量. 证明此有限元定义良好. 使用该元, 对具有 Neumann 边界条件的 $\Delta u=0$ 导出第~\ref{sec:ch05-poisson-variational-approximation} 节结果的类似结论, 即定理~\ref{thm:ch05-poisson-h1-error} 中取 $m=4$. 提示: 先证明有限元 $(K,P_3,\widetilde N)$ 定义良好, 其中
\[
\MathBookFormulaID{formula-ch05-exercise-05-augmented-degrees-of-freedom}
\widetilde N=N\cup\left\{v\longmapsto
\int_{\partial K}\frac{\partial v}{\partial\nu}\,ds\right\},
\]
再证明 $u-\widetilde I^h u$ 的估计, 并注意 $\Delta u=0$ 蕴含
\[
\MathBookFormulaID{formula-ch05-exercise-05-harmonic-element-flux}
\int_{\partial K}\frac{\partial u}{\partial\nu}\,ds=0
\quad\forall K.
\]
\end{MBExercise}

\begin{MBExercise}
\label{exr:ch05-06}
证明例~\ref{ex:ch05-mixed-boundary-singularity} 中使用的反射把空间
\[
\MathBookFormulaID{formula-ch05-exercise-06-reflected-space}
\left\{v\in H^2(\Omega):v|_\Gamma=0,
\left.\frac{\partial v}{\partial\nu}\right|_{\partial\Omega\setminus\Gamma}=0\right\}
\]
映到 $H^2(\widetilde\Omega)$.
\end{MBExercise}

\begin{MBExercise}
\label{exr:ch05-07}
证明变分问题 $a(u,v)=(f,v)$ 对所有 $v\in V$ 成立, 其中形式为~\eqref{eq:ch05-general-nonsymmetric-form}, 且 $V=H^1(\Omega)$, 等价于 $Au=f$, 其中 $A$ 由式~\eqref{eq:ch05-general-nonsymmetric-operator} 定义, 边界条件为~\eqref{eq:ch05-general-natural-boundary-condition}.
\end{MBExercise}

\begin{MBExercise}
\label{exr:ch05-08}
对具有形式~\eqref{eq:ch05-general-nonsymmetric-form} 的一般二阶椭圆 Dirichlet 问题表述对偶性论证. 提示: 考虑伴随算子的 Dirichlet 问题
\[
\MathBookFormulaID{formula-ch05-exercise-08-adjoint-operator}
A^tv:=-\sum_{i,j=1}^n\frac{\partial}{\partial x_i}
\left(a_{ij}\frac{\partial v}{\partial x_j}\right)
-\sum_{k=1}^n\frac{\partial}{\partial x_k}(b_kv)+b_0v.
\]
陈述所涉及的全部椭圆问题所需的正则性.
\end{MBExercise}

\begin{MBExercise}
\label{exr:ch05-09}
确定练习~\ref{exr:ch05-08} 中伴随问题相应的自然边界条件. 对具有形式~\eqref{eq:ch05-general-nonsymmetric-form} 和空间 $V=H^1(\Omega)$ 的一般二阶椭圆 Neumann 问题表述对偶性论证.
\end{MBExercise}

\begin{MBExercise}
\label{exr:ch05-10}
令 $u_h$ 为式~\eqref{eq:ch05-discrete-inhomogeneous-poisson} 的解. 证明
\MathBookSourcePage{165}
\[
\MathBookFormulaID{formula-ch05-exercise-10-inhomogeneous-h1-error}
\|u-u_h\|_{H^1(\Omega)}
\leq\inf_{v\in V_h}\|u-g_D-v\|_{H^1(\Omega)}
+2\|g_D-I^h g_D\|_{H^1(\Omega)}.
\]
假设椭圆正则性~\eqref{eq:ch05-elliptic-regularity-h2} 成立, 再证明
\[
\MathBookFormulaID{formula-ch05-exercise-10-inhomogeneous-l2-error}
\|u-u_h\|_{L^2(\Omega)}
\leq C\left[h\inf_{v\in V_h}\|u-g_D-v\|_{H^1(\Omega)}
+\|g_D-I^h g_D\|_{L^2(\Omega)}\right].
\]
\end{MBExercise}

\begin{MBExercise}
\label{exr:ch05-11}
证明当式~\eqref{eq:ch05-discrete-inhomogeneous-poisson} 中的 $f,g_N$ 分别由其插值函数代替时, 所得扰动问题的解满足定理~\ref{thm:ch05-poisson-h1-error} 和~\ref{thm:ch05-poisson-l2-error} 的类似结论.
\end{MBExercise}

\begin{MBExercise}
\label{exr:ch05-12}
考虑方程~\eqref{eq:ch05-poisson-equation}, 其中 $u=0$ 于 $\Gamma$, 并在 $\partial\Omega\setminus\Gamma$ 上取 Robin 边界条件
\[
\MathBookFormulaID{formula-ch05-exercise-12-robin-condition}
\alpha u+\frac{\partial u}{\partial\nu}=0
\quad\text{于 }\partial\Omega\setminus\Gamma,
\]
其中 $\alpha$ 为常数.
\begin{enumerate}
\item[(1)] 导出该问题的变分表述, 并给出保证它在 $H^1$ 中有唯一解的 $\alpha$ 和 $f$ 的条件(若需要). 提示: 乘以 $v$ 并分部积分, 使用边界条件把 $(\partial u/\partial\nu)v$ 项转为 $uv$ 项.
\item[(2)] 证明若 $u\in H^2$ 满足你的变分问题, 且 $f\in L^2$, 则它满足本练习中的原微分问题.
\item[(3)] 使用一般网格上的分片线性函数, 在 $H^1$ 和 $L^2$ 中证明变分逼近的误差估计. 假设问题具有 $H^2$ 正则性.
\item[(4)] 对 $\Omega=[0,1]\times[0,1]$ 上的正则网格使用分片线性函数, 计算变分逼近相应的``差分模板''(参见第~\ref{sec:ch00-relation-to-difference-methods} 节). 提示: 若网格由 $45^\circ$ 直角三角形组成, 则远离边界处为标准五点差分模板.
\end{enumerate}
\end{MBExercise}

\begin{MBExercise}
\label{exr:ch05-13}
证明对 $p\geq1$,
\[
\MathBookFormulaID{formula-ch05-exercise-13-lp-boundary-poincare}
\|v\|_{L^p(\Omega)}
\leq C\left(\left|\int_\Gamma v\,ds\right|+|v|_{W_p^1(\Omega)}\right)
\quad\forall v\in W_p^1(\Omega),
\]
其中 $\Gamma\subset\partial\Omega$ 具有正测度, 常数 $C$ 只依赖于 $\Omega$ 和 $\Gamma$.
\end{MBExercise}

\begin{MBExercise}
\label{exr:ch05-14}
考虑 $\Omega=[0,1]\times[0,1]$ 上具有周期边界条件的方程~\eqref{eq:ch05-poisson-equation}. 证明其变分表述与 Neumann 问题相似, 只是变分空间中要施加周期性.
\end{MBExercise}

\begin{MBExercise}
\label{exr:ch05-15}
对双线性形式~\eqref{eq:ch05-symmetric-elliptic-bilinear-form} 证明命题~\ref{prop:ch05-mixed-boundary-coercivity}.
\end{MBExercise}

\begin{MBExercise}
\label{exr:ch05-16}
在条件~\eqref{eq:ch05-general-form-continuity}、\eqref{eq:ch05-shifted-coercivity} 以及连续问题唯一可解并满足椭圆正则性估计~\eqref{eq:ch05-general-h2-regularity} 的假设下, 证明不定椭圆问题的离散问题在 $h$ 足够小时唯一可解.
\end{MBExercise}

\MathBookSourcePage{166}
\begin{MBExercise}
\label{exr:ch05-17}
证明定理~\ref{thm:ch05-plate-negative-norm-error}. 提示: 沿用第~\ref{sec:ch05-negative-norm-estimates} 节的技术.
\end{MBExercise}

\begin{MBExercise}
\label{exr:ch05-18}
证明次数小于 $N$、在某个 Sobolev 范数中闭且有界的多项式集合是紧集. 提示: 把该范数写成多项式系数上的度量, 说明该集合对应于 Euclidean 空间中的闭有界子集.
\end{MBExercise}

\begin{MBExercise}
\label{exr:ch05-19}
为板弯曲问题表述非齐次边值问题.
\end{MBExercise}

\begin{MBExercise}
\label{exr:ch05-20}
使用式~\eqref{eq:ch05-plate-garding-inequality} 对 $-3<\nu<1$ 证明式~\eqref{eq:ch05-plate-coercivity-computation}. 提示: 使用 Rellich 引理(Agmon 1965), 它断言 $H^m(\Omega)$ 的有界子集在 $L^2(\Omega)$ 中预紧, 其中 $m\geq1$.
\end{MBExercise}

\begin{MBExercise}
\label{exr:ch05-21}
设 $T_h=\{(r,\theta):0<r<h,\ \theta_0<\theta<\theta_1\}$. 证明
\[
\MathBookFormulaID{formula-ch05-exercise-21-sector-scaling}
\inf_{\widetilde v\in P_k^2}\int_{T_h}
|\nabla(r^\beta\sin\beta\theta)-\widetilde v|^2\,dx
=c_{k,\beta}h^{2\beta},
\]
其中 $P_k$ 表示次数至多为 $k$ 的多项式, $c_{k,\beta}>0$ 只依赖于 $k,\beta$. 提示: 作坐标伸缩, 证明上式成立且
\[
\MathBookFormulaID{formula-ch05-exercise-21-unit-sector-constant}
c_{k,\beta}=\inf_{\widetilde v\in P_k^2}\int_{T_1}
|\nabla(r^\beta\sin\beta\theta)-\widetilde v|^2\,dx.
\]
再用紧性论证证明 $c_{k,\beta}=0$ 将导致矛盾 $\nabla(r^\beta\sin\beta\theta)\in P_k^2$.
\end{MBExercise}

\begin{MBExercise}
\label{exr:ch05-22}
证明对任意 $f\in L^p(\Omega)$ 和任意 $\widetilde\Omega\subset\Omega$,
\[
\MathBookFormulaID{formula-ch05-exercise-22-restricted-domain-best-error}
\inf_{v\in W}\int_\Omega|f-v|^p\,dx
\geq\inf_{v\in W}\int_{\widetilde\Omega}|f-v|^p\,dx
\]
对任意子集 $W\subset L^p(\Omega)$ 成立.
\end{MBExercise}

\begin{MBExercise}
\label{exr:ch05-23}
设 $\Omega$ 为单位圆盘, $\Omega_h$ 为内接于 $\Omega$、边数为 $1/h$ 的正多边形序列. 在每个 $\Omega_h$ 上考虑简支板问题. 证明这些问题的解在 $h\to0$ 时收敛, 但不收敛到 $\Omega$ 上简支板问题的解. 这称为 Babuška 悖论. 保证此情形下有限元逼近收敛到正确解的技术见 Scott (1977a).
\end{MBExercise}
